\section*{Chapter \ref{ch:rs:vectorised-backtests} --- Vectorised Backtests}

\begin{solution}{exo:rs:vectorised-backtests:1}
The gross return is $12\%/252 = 0.0476\%$ a day; by \cref{prop:rs:vectorised-backtests:breakeven} the break-even cost is
$0.0476\%/(2 \times 0.2) = 0.119\%$, 11.9 basis points per unit traded.
\end{solution}

\begin{solution}{exo:rs:vectorised-backtests:2}
After the period A is worth $0.66$ and B $0.40$ of a capital of 1.06: weights 62.3\% and 37.7\%. To return to 60/40, sell 2.3\% of capital of A and buy
2.3\% of B.
\end{solution}

\begin{solution}{exo:rs:vectorised-backtests:3}
Compounded: $(1.5 \times 0.6)^5 = 0.9^5 = 0.59$, a loss of 41\%. Summed: $5 \times 50\% - 5 \times 40\% = +50\%$, a gain that never happened.
\end{solution}

\begin{solution}{exo:rs:vectorised-backtests:4}
The reversal's signal is minus the day's return, so multiplying it by that return gives minus the return squared: a certain loss. A momentum signal that
includes today's return gives plus a positive multiple of it squared: a certain gain. Look-ahead adds the signal's correlation with the return it was
built from, whatever its sign.
\end{solution}

\begin{solution}{exo:rs:vectorised-backtests:5}
A survivors' universe removes the names that later delisted, and with them their returns before and at delisting. A strategy that would have held those
names long loses the losses (its \omterm{def:rs:vectorised-backtests:backtest}{backtest} rises); one that would have held them short loses the gains (its \omterm{def:rs:vectorised-backtests:backtest}{backtest} falls). In the synthetic market the
delisted names were momentum's past losers, held short, and their falls to delisting were profits for the short side.
\end{solution}

\begin{solution}{exo:rs:vectorised-backtests:6}
A dollar-neutral book has net exposure 0, so its cash is its capital: $+3\%$ a year of interest. Its shorts are 1 (half of gross 2) and cost $0.5\% \times 1
= 0.5\%$. Net: $+2.5\%$ a year. (The Sharpe ratio is computed on the return in excess of the cash rate.)
\end{solution}

\begin{solution}{exo:rs:vectorised-backtests:7}
\texttt{rs\_vecbt.smoothed(3)}: turnover falls from 75\% to 42\% of the book a day and the gross Sharpe ratio from 6.95 to 4.13, while the break-even cost
barely moves (7.4 basis points against 7.3): the smoothing removes alpha along with trading, because the reversal's information is one day old. It helps
only a signal whose information outlasts its noise.
\end{solution}

\begin{solution}{exo:rs:vectorised-backtests:8}
Market-on-close orders are submitted before the close, typically minutes before, when the closing price is not known; a signal computed from that price is
not available when the order must be sent. Either compute the signal from prices known before the submission cutoff, or trade at the next close; and model
the auction's own costs.
\end{solution}

\begin{solution}{pb:rs:vectorised-backtests:1}
\begin{enumerate}
\item Sharpe ratio 7.17, 27.0\% a year, assuming it trades at the close it computed its signal from, in today's survivors, without costs.
\item $-65$, with $-694\%$ a year: the signal is minus the day's return, multiplied by that return.
\item 74\% of the book a day.
\item The returns up to the close, but not the ability to trade at that close, nor the list of survivors.
\item Reversal: 7.22 on the point-in-time universe, 6.95 on the 500 most liquid. Momentum: 0.44, then 0.32.
\item The names that delisted were momentum's losers, held short; dropping them drops the short side's profits (0.21 against 0.44).
\item Raise it: distressed stocks that failed would be missing from the \omterm{def:rs:vectorised-backtests:backtest}{backtest}'s longs.
\item Reversal: 6.95 to $-2.58$. Momentum: 0.32 to 0.21.
\item 7.3 basis points per unit traded.
\item About 0.25\% a year for the reversal: $-2.58$ to $-2.64$.
\item $-9.45$: the reversal's information lasts one day, so a book traded a day late earns no alpha and still pays its costs.
\item Compounded, capital falls to 0.35 of its start; summed, it reaches $-0.04$.
\item \textbf{Named result.} Reversal: 7.17 (naive), 7.22 (point-in-time universe), 6.95 (liquid), $-2.58$ (costs), $-2.64$ (borrow), $-9.45$ (next close).
Momentum: 0.21, 0.44, 0.32, 0.21, 0.19, 0.17. The reversal dies at the costs and again at the timing; momentum loses 60\% of its best number but keeps a
positive, small Sharpe ratio through every correction.
\item Momentum is a level-1 result: slow, liquid, insensitive to execution. The reversal needs the closing auction modelled (level~2 with an auction fill
model, or level~3), and its level-1 answer is already negative.
\item It would give the reversal a legitimate way to trade at the close; whether any alpha is left depends on how much of the day's move happens after
15:50, and the auction's cost replaces the generic 10 basis points.
\item Trade only names whose signal crosses a threshold, keep positions until the signal reverses (a buy/hold spread, as Novy-Marx and Velikov recommend),
or combine the signal with slower ones; each cuts turnover more than alpha only if the alpha is not all in the first day.
\item The first two lines: the \omterm{def:rs:vectorised-backtests:timing}{execution lag} and the point-in-time universe.
\item That a daily reversal, with turnover far above 50\% a month, does not survive costs unless its costs are exceptionally low.
\item Every variant tried, with its settings (universe, lag, costs), the level-1 waterfall, and the decision taken (chapter~1).
\item Screening ideas quickly and honestly when execution does not change the answer.
\end{enumerate}
\end{solution}

\begin{solution}{iq:rs:vectorised-backtests:1}
Look-ahead (inputs not known at the \omterm{def:rs:vectorised-backtests:timing}{decision time}, or trading at the price the signal used), survivorship (today's universe), no or low costs, no borrow or
financing, and data snooping (many variants, one reported); also arithmetic cumulation of large returns, and ignoring capacity.
\end{solution}

\begin{solution}{iq:rs:vectorised-backtests:2}
Build the signal panel with each input as of its \omterm{def:rs:point-in-time-data-and-the-biases:bitemporal}{knowledge time}; mask to the point-in-time \omterm{def:rs:universe-symbology-and-corporate-actions:universe}{tradable universe}; convert to target weights; lag them to the
first period they can be held; compute the drifted book and the trades to the new target; charge costs on trades, borrow on shorts, financing on cash and
leverage; compound. Biases enter at each step: timing (lag), universe (survivors), costs, financing, compounding.
\end{solution}

\begin{solution}{iq:rs:vectorised-backtests:3}
Its turnover may make costs exceed its gross return (the break-even cost may be below realistic costs); it may depend on trading at a price it could not
get; it may not scale (capacity); it may be one of many variants tried. On the synthetic reversal, 7 became $-2.6$ at 10 basis points of cost.
\end{solution}

\begin{solution}{iq:rs:vectorised-backtests:4}
Good enough for slow, liquid strategies whose trades are small against volume and costs are a spread per unit traded. An event-driven (or replay) \omterm{def:rs:vectorised-backtests:backtest}{backtest} is
needed when fills depend on the order book (passive or limit orders), when timing within the day or auctions matter, when the strategy's own impact matters,
or when events intervene between decision and trade.
\end{solution}

\begin{solution}{iq:rs:vectorised-backtests:5}
Charge borrow fees by name on the short positions held each period (with availability limits), credit interest on cash and short proceeds as the prime broker
pays it, charge the financing spread on borrowed cash for leverage, and compute the Sharpe ratio on the return in excess of the cash rate.
\end{solution}

\begin{solution}{iq:rs:vectorised-backtests:6}
Per period the book trades $2\tau$ of weight and pays $2\tau c$; the net return is $g - 2\tau c$, zero at $c = g/(2\tau)$. The synthetic reversal: $g$ about
0.109\% a day and $\tau = 0.75$, so 7.3 basis points; monthly momentum: $g$ about 0.011\% a day and $\tau$ about 0.019, so about 30 basis points: four times the
reversal's tolerance for cost, with a tenth of its gross return.
\end{solution}
